\documentclass[a4paper,11pt,fleqn]{article}%fleqn pour aligner les equations à gauche
\usepackage{graphicx}
\usepackage[utf8]{inputenc}  
\usepackage[T1]{fontenc}
\usepackage{lscape}

\usepackage{tgheros,textcomp}%Police Arial
\renewcommand{\familydefault}{\sfdefault}

%\setlength{\mathindent}{0pt}%indentation dans math = 0

\usepackage[top = 1cm, bottom = 2cm, left = 1cm, right = 1cm]{geometry}
\usepackage{titlesec}
%\setcounter{secnumdepth}{4}
\usepackage{xcolor}
\usepackage{amssymb} %double flèche de réaction
\usepackage{mathrsfs}%farraday
\usepackage{xtab}%tableau sur plusieurs pages
\usepackage{esvect}%grands vecteurs

\usepackage[version=4]{mhchem} %equations chimiques

\usepackage{array}
\usepackage{chemfig}
\usepackage{multirow}
%\usepackage{sectsty}

\usepackage{caption}
\usepackage{adjustbox}
\usepackage[labelformat=simple]{subfig}
\usepackage{float}


\usepackage{amsmath}
\usepackage{amssymb}
\usepackage[cal=boondoxo]{mathalfa} 
\usepackage{enumitem}
\usepackage{lastpage}
\usepackage{tcolorbox}
\usepackage{siunitx}
\usepackage[european, straightvoltages, RPvoltages]{circuitikz}
\usepackage{tikz}
\usetikzlibrary{shapes.geometric}
\usetikzlibrary{calc}
\usetikzlibrary{automata,positioning}
\usepackage{multicol}
\usepackage{eqnarray}
\usepackage{tabularx,colortbl}%colorer fond cellules

\sisetup{
	detect-all,
	output-decimal-marker={,},
	group-minimum-digits = 3,
	group-separator={~},
	number-unit-separator={~},
	inter-unit-product={~}
} %setup pour SIunitx

\usepackage{listings}%pour insérer le code python
\definecolor{codegreen}{rgb}{0,0.6,0}
\definecolor{codegray}{rgb}{0.5,0.5,0.5}
\definecolor{codepurple}{rgb}{0.58,0,0.82}
\definecolor{backcolour}{rgb}{0.95,0.95,0.92}


%\graphicspath{ {/home/rweye/Documents/Enseignement/2023/Cours/TSPE/Eval/Ressources/} }

\titleformat{\section}
{\titlerule
\vspace{.8ex}%
\normalfont \bfseries}
{\thesection :}{.5em}{\bfseries}

\renewcommand{\thesection}{Exercice \arabic{section}}
\titlespacing{\section}{0pt}{1em}{1em}
\renewcommand{\thesubsection}{\hspace{1cm}\alph{subsection}.}
\renewcommand\thesubfigure{\textbf{Document \arabic{subfigure} :}}

\title{\vspace{-1cm}\large 
	\begin{tabular}{|m{5cm}|>{\centering}m{8cm}|>{\centering\arraybackslash}m{4cm}|}
		\hline
		Nom : & \textbf{\textsc{DS TSPE N°3}} &  Note : \\
		Prénom : & & $\_\_\_$ \textbf{20} \\ \hline
\end{tabular}}
\date{}

\newpagestyle{mystyle}{%
	\setfoot{ROSALIE}{\thepage\,$|$\,\pageref{LastPage}}{}
}
\renewpagestyle{plain}{%
	\setfoot{ROSALIE}{\thepage\,$|$\,\pageref{LastPage}}{}
}

\pagestyle{mystyle}

\newenvironment{itemize*}%
{\begin{itemize}[label=$-$]%
		\setlength{\itemsep}{0pt}%
		\setlength{\parskip}{0pt}}%
	{\end{itemize}}

\newenvironment{itemize**}%
{\begin{itemize}[label=$\_\_$,labelindent=0pt, leftmargin=*,topsep=0pt]%
		\setlength{\itemsep}{0pt}%
		\setlength{\parskip}{0pt}}%
	{\end{itemize}}

\renewcommand{\arraystretch}{1.5}


\begin{document}
\maketitle
\vspace{-2.5cm}
\section{Casque à réduction active de bruit}
La réduction active de bruit consiste à supprimer le bruit résiduel dans les oreillettes per émission d’un signal appelé « anti-bruit ». Le Document 1 présente le niveau d'intensité sonore mesuré en fonction de la fréquence dans les situations suivantes :
\begin{itemize*}
	\item  sans casque (cas n°1)
	\item avec casque en réduction passive (cas n°2)
	\item avec casque en réduction passive et active (cas n°3)
\end{itemize*}
\begin{enumerate}
	\item Identifier approximativement les domaines de fréquence pour lesquels :
	\begin{itemize*}
		\item seule la réduction passive est efficace
		\item seule la réduction active est efficace
		\item les deux réductions sont efficaces
	\end{itemize*}
\end{enumerate}
On mesure l'intensité sonore lorsque le bruit et l'anti-bruit sont émis simultanément dans trois situations (Document 2). Le niveau sonore du bruit et de l'anti-bruit lorsqu'ils sont émis seuls est de \SI{50}{dB}.
\begin{enumerate}[resume]
	\item \begin{enumerate}
		\item Préciser dans quelle situation l'intensité sonore est la somme des intensités des sons pris séparément.
		\item Déterminer la situation qu'il faut retenir pour le dispositif de réduction active de bruit.
		\item Calculer l'atténuation associée.
	\end{enumerate}
\end{enumerate}

\begin{center}
\includegraphics[scale=.55]{Ressources/TSPE_DS3b_E1_doc1.png}
\end{center}
\textbf{Document 2 : Variation de l'intensité sonore}\\

\begin{tabular}{|>{\centering}m{4cm}|>{\centering}m{4cm}|>{\centering}m{4cm}|>{\centering \arraybackslash}m{4cm}|} \hline
\textbf{Situation} & a & b & c\\ \hline
\textbf{Verrouillage de phase} & Non & En phase & En opposition de phase \\ \hline
\textbf{Niveau sonore (\unit{\deci\bel})} & \num{53} $\pm$ \num{1} & \num{56} $\pm$ \num{1} &\num{44} $\pm$ \num{1} \\ \hline
\end{tabular}

\newpage
\section{Tube cathodique}
\emph{Les deux parties de l'exercice peuvent être traitées indépendamment}\\
\textbf{Partie 1 : Accélération}\\
Avant le développement et la généralisation des écrans LCD, les oscilloscopes, les écrans d'ordinateurs et les téléviseurs étaient constitués d'un tube cathodique.\\
A la base du tube cathodique, un canon à électrons émet et accélère des électrons en direction de l'écran. (Document 1).\\ Lors d'une première phase du trajet, les électrons passent entre deux plaques verticales percées. Lors de cette phase, les électrons passent d'une vitesse négligeable à une vitesse notée $v_0$ par la suite.
\begin{enumerate}
	\item Les deux plaques sont séparées d'une distance $d$ = \SI{5,0}{cm} et la tension aux bornes des plaques est de \SI{200}{V}. Montrer que le  champ électrique régnant entre les deux plaques a une valeur absolue de \SI{4.0E3}{\volt\per\meter}.
	\item Montrer que le poids de l'électron est négligeable par rapport à la force électrique qu'il subit.
	\item Calculer le travail de la force électrique appliquée sur la distance $d$.
	\item Utiliser le théorème de l'énergie cinétique pour déterminer la valeur de la vitesse de l'électron à la sortie des plaques.
\end{enumerate}

\textbf{Partie 2 : Déviation}\\
Suite à cette accélération, les électrons passent à l'intérieur de deux condensateurs plans $C_1$ et $C_2$ chargés qui permettent de dévier les électrons horizontalement ($C_1$) et verticalement ($C_2$).\\
En fin de course, les électrons impactent l'écran sur lequel est déposée une couche électroluminescente.\\
On étudie la déflexion dans le condensateur $C_2$.
\begin{enumerate}[resume]
	\item L'électron pénètre dans le condensateur $C_2$ avec une vitesse horizontale $\vec{v_0}$ (Document 2). Représenter la force électrique $\vec{F_e}$ subie.
\end{enumerate} 
En considérant que $\vec{F_e}$, les équations horaires sont : 
$$\left\lbrace \begin{array}{l}
	x(t) = v_0 \cdot t \\
	y(t) = - \frac{1}{2}\cdot \dfrac{e\cdot E}{m} \cdot t^2
	\end{array} \right \rbrace$$
\begin{enumerate}[resume]
	\item Déterminer l'équation de sa trajectoire et sa nature.
	\item Préciser le mouvement de l'électron après sa sortie du condensateur.
\end{enumerate}

\includegraphics[scale=1.2]{Ressources/TSPE_DS3c_E2_doc1.png} 
\includegraphics[scale=1.2]{Ressources/TSPE_DS3c_E2_doc2.png}

\textbf{Données :}
\begin{itemize}
	\item Masse d'un électron : $m_{\ce{e-}}$ = \SI{9,11E-31}{kg}
	\item Intensité de pesanteur : $g$ = \SI{9,81}{\newton\per\kilo\gram}
	\item Charge élémentaire : $e$ = \SI{1,60 E-19}{C}
\end{itemize}
\end{document}